\label{intro_md_docs_intro}%
\Hypertarget{intro_md_docs_intro}%
 The Hankel transform, also known as the Fourier–\+Bessel transform, is widely used in scattering. The transform is needed in many areas, e.\+g. to calculate scattering functions, to correct for multiple scattering, to simulate spin echo signals like those obtained on Larmor at the ISIS Neutron and Muon Source, as well as to qualitatively analyse phase contrast images (PCI) taken at the Diamond synchrotron source.

Small-\/angle scattering (SAS) is an experimental technique using X-\/rays (SAXS) or neutrons (SANS) to study the structure of materials on the nanoscale (roughly 1–1000 nm). SAS measures how radiation is scattered at very small angles, which tells us about large-\/scale structures inside the sample.

In this library, there are currently ... different implementations of the Hankel transform, correspnding to different ways of approximating the Hankel transform integral.

\tabulinesep=1mm
\begin{longtabu}spread 0pt [c]{*{3}{|X[-1]}|}
\hline
\cellcolor{\tableheadbgcolor}\textbf{ index   }&\PBS\centering \cellcolor{\tableheadbgcolor}\textbf{ Strategy Name   }&\PBS\centering \cellcolor{\tableheadbgcolor}\textbf{ Type    }\\\cline{1-3}
\endfirsthead
\hline
\endfoot
\hline
\cellcolor{\tableheadbgcolor}\textbf{ index   }&\PBS\centering \cellcolor{\tableheadbgcolor}\textbf{ Strategy Name   }&\PBS\centering \cellcolor{\tableheadbgcolor}\textbf{ Type    }\\\cline{1-3}
\endhead
0   &\PBS\centering DE   &Double-\/exponential quadrature    \\\cline{1-3}
1   &\PBS\centering DE   &Double-\/exponential quadrature    \\\cline{1-3}
12   &\PBS\centering QWE\+\_\+\+Key   &QWE with continued fraction expansion    \\\cline{1-3}
13   &\PBS\centering QWE\+\_\+\+Chave   &QWE with Shanks transformation    \\\cline{1-3}
...   &\PBS\centering ...   &...   \\\cline{1-3}
\end{longtabu}
